\documentclass[12pt,a4paper]{article}
\usepackage{amsmath,amssymb,amsfonts,booktabs,array,microtype,xcolor,graphicx,hyperref}
\usepackage{amsthm}

\newcommand{\PSL}{\mathrm{PSL}(2,\mathbb{C})}
\newcommand{\mfld}[1]{\texttt{#1}}
\newcommand{\word}[1]{\texttt{#1}}
\newcommand{\degr}{^{\circ}}

\newtheorem{conjecture}{Conjecture}
\newtheorem{proposition}{Proposition}
\newtheorem{remark}{Remark}

\begin{document}

\title{Homology Class Asymmetry in the Loxodromic Twist Spectrum\\
of Compact Hyperbolic 3-Manifolds}

\author{Marvin L.~Gentry\thanks{Independent Researcher, Seattle, WA, USA.
  \texttt{drmlgentry@protonmail.com}. ORCID: 0009-0006-4550-2663.}}
\date{\today}
\maketitle

\begin{abstract}
We introduce a homology-resolved extreme-value invariant of the complex length
spectrum of a compact hyperbolic 3-manifold.
For manifolds with $H_1(M,\mathbb{Z})=\mathbb{Z}/5$, the abelianization map
partitions loxodromic elements into five homology classes.
We study the spectral floor $\phi_{\rm floor}^{(k)}(L)$---the minimum folded
twist angle among loxodromic words of length $\leq L$ in class $k$---for the
manifold \mfld{m006} (OrientableClosedCensus[43], $\mathrm{vol}=2.029$,
$H_1=\mathbb{Z}/5$).
A complete census up to word length~12 ($1{,}062{,}802$ genuine loxodromics)
shows that all five floors decrease rapidly with $L$ and are well described by
exponential decay over the observed range, with class-dependent rates
$c_k$ forming a paired hierarchy:
$c_1=c_4\approx 0.985$, $c_2=c_3\approx 0.401$, $c_0\approx 0.750$,
confirmed by Akaike Information Criterion (AIC) model selection against
power-law alternatives ($\Delta\mathrm{AIC}>40$ for classes~0,1,4).
At $L=12$ the floors satisfy an approximate $4:2:1$ ratio.
The same pairing structure appears for the Meyerhoff manifold \mfld{m003}
(also $H_1=\mathbb{Z}/5$), suggesting the hierarchy is a property of
$\mathbb{Z}/5$ torsion rather than of either specific manifold.
We conjecture that the pairing $c_k=c_{5-k}$ reflects an arithmetic symmetry
of the holonomy representation, arising from complex conjugation in the
invariant trace field.
The floor $\phi_{\rm floor}^{(k)}(L)\to 0$ for all $k$ follows from the
Ambient Prime Geodesic Theorem~\cite{BalogBiggs23}; no claim is made
regarding the asymptotic form of the decay rate.
\end{abstract}

\section{Introduction}
\label{sec:intro}

The complex length spectrum $\{\ell(\gamma)+i\phi(\gamma)\}_\gamma$ of a
compact hyperbolic 3-manifold $M$ is a fundamental geometric invariant,
encoding geodesic lengths $\ell(\gamma)$ and loxodromic twist angles
$\phi(\gamma)=\mathrm{Im}\log\lambda(\gamma)$.
While the length spectrum has been extensively studied~\cite{MaclachlanReid},
the distribution of twist angles and their relationship to the homological
structure of $M$ is less well understood.

In this paper we study the \emph{spectral floor}
$\phi_{\rm floor}^{(k)}(L)=\min_{\gamma\in C_k,\,|w(\gamma)|\leq L}
|\phi_{\rm fold}(\gamma)|$
for each homology class $k\in H_1(M,\mathbb{Z})=\mathbb{Z}/5$,
as a function of word length~$L$.
Here $\phi_{\rm fold}=\min(|\phi|\bmod 180\degr,\;180\degr-|\phi|\bmod 180\degr)$
folds the twist to $[0\degr,90\degr]$, and $C_k$ denotes the set of
loxodromic elements in homology class~$k$.

\textbf{Main result.}
For the manifold \mfld{m006} (OrientableClosedCensus[43], $\mathrm{vol}=2.029$,
$H_1=\mathbb{Z}/5$), a census of all words up to length~12 reveals that
$\phi_{\rm floor}^{(k)}(L)$ is well described by exponential decay over the observed range for all five
homology classes, but with class-dependent rates forming a $4:2:1$ hierarchy:
\begin{equation}
  \frac{\phi^{(0)}}{\phi^{(2,3)}} \approx 2,\qquad
  \frac{\phi^{(2,3)}}{\phi^{(1,4)}} \approx 2,
  \label{eq:ratio}
\end{equation}
stable at word length~12 and confirmed over lengths 1--12.
The paired structure $c_1=c_4$ and $c_2=c_3$ is consistent with an
outer involution of $\pi_1(\mfld{m006})$ exchanging classes
$k\leftrightarrow 5-k$ (mod~5).

\textit{Motivation.}
This study was originally motivated by the observation, made in a
companion series of papers~\cite{GentryCKM,GentryPMNS,GentryTwist},
that \mfld{m003} and \mfld{m006} reproduce Standard Model quark and
lepton mixing matrices with no free parameters.
The present paper is independent of that program: the spectral floor
invariant $\phi_{\rm floor}^{(k)}(L)$ and the decay rate hierarchy are
defined and studied purely as geometric objects, and all results stand
independently of any physical interpretation.

\section{Setup}
\label{sec:setup}

\subsection{Manifolds}
\begin{itemize}
\item $M_1=\mfld{m003}$ (OrientableClosedCensus[1], Meyerhoff manifold~\cite{Meyerhoff87},
      $\mathrm{vol}=0.9814$, $H_1=\mathbb{Z}/5$).
\item $M_2=\mfld{m006}$ (OrientableClosedCensus[43], $\mathrm{vol}=2.029$,
      $H_1=\mathbb{Z}/5$).
\end{itemize}
Both manifolds are likely arithmetic~\cite{MaclachlanReid,Chinburg98};
\mfld{m003} is proven arithmetic with invariant trace field of
discriminant~$-283$~\cite{Chinburg98}.

\subsection{Homology class}
The abelianization $\pi_1(M)\to H_1(M,\mathbb{Z})=\mathbb{Z}/5$ assigns to
each word $w=g_1\cdots g_n$ the class
\begin{equation}
  [w] = \sum_{i=1}^n \varepsilon(g_i) \bmod 5,\qquad
  \varepsilon(g) = \begin{cases}+1 & g\text{ lowercase}\\-1 & g\text{ uppercase}\end{cases}.
\end{equation}

\subsection{Census procedure}
We enumerate all reduced words up to length~$L$ over the generators
$\{a,b,A,B\}$ of $\pi_1(\mfld{m006})$ (where $A=a^{-1}$, $B=b^{-1}$),
compute the holonomy matrix $\rho(w)\in\PSL$ via SnapPy's
\texttt{G.SL2C(w)} method~\cite{SnapPy}, and record $\phi_{\rm fold}(w)$
for loxodromic elements ($|\mathrm{Tr}(\rho(w))|>2.01$).
At length~12 this yields $1{,}062{,}802$ genuine loxodromics from
$1{,}062{,}880$ words evaluated in $84.6$~seconds.
Twist angles are computed as $\phi=\arg\lambda$ where $\lambda$ is the
eigenvalue of larger modulus of the $2\times 2$ holonomy matrix.
Near-$\pi$ twists ($\phi_{\rm fold}<0.01\degr$) are verified by computing
$|\mathrm{Im}(\lambda)/\mathrm{Re}(\lambda)|$ directly; all values
reported are stable to the numerical precision of \textsc{SnapPy}'s
default 15-digit arithmetic.
Words with $|\lambda|\leq 1.01$ (mapping to $\pm I$ in $\PSL$) are excluded
as non-geodesic relators.

\section{Results}
\label{sec:results}

\subsection{Spectral floors at length 7: initial asymmetry}
Table~\ref{tab:floors_L7} shows the length-7 floors for both manifolds.
On \mfld{m006}, class~4 already has a floor of $0.005\degr$ (word
\word{AAABAB}), while the other four classes have floors $\geq 1.6\degr$
--- a $637\times$ asymmetry.
On \mfld{m003}, the floors are within $2\times$ of each other.
This contrast between the two manifolds at length~7 motivates
extending the census to longer lengths for both manifolds.

\begin{table}[tp]
\caption{Spectral floors per homology class at word length $\leq 7$
  (genuine loxodromics only, $|\lambda|>1.01$).}
\label{tab:floors_L7}
\begin{tabular}{cllll}
\toprule
Class $k$ & \multicolumn{2}{c}{\mfld{m003}} &
            \multicolumn{2}{c}{\mfld{m006}} \\
          & $\phi_{\rm floor}$ ($\degr$) & Word
          & $\phi_{\rm floor}$ ($\degr$) & Word \\
\midrule
0 & 2.643 & \word{AAAAB}   & 2.132 & \word{AAABB}   \\
1 & 2.763 & \word{AAAbbbb} & 3.374 & \word{BBBB}    \\
2 & 4.857 & \word{AABBAbb} & 2.399 & \word{AABABab} \\
3 & 5.376 & \word{AABABBB} & 1.611 & \word{ABaBAb}  \\
4 & 3.269 & \word{AAb}     & \textbf{0.005} & \word{AAABAB}  \\
\midrule
\multicolumn{2}{l}{Max/min ratio} & $2\times$
& & $637\times$ \\
\bottomrule
\end{tabular}
\end{table}

\subsection{Length-12 floors for both manifolds}

Table~\ref{tab:floors_both_L12} compares length-12 floors for both
manifolds.  The $4:2:1$ pairing structure appears in both cases.

\begin{table}[tp]
\caption{Spectral floors per homology class at word length $\leq 12$
  for both manifolds with $H_1=\mathbb{Z}/5$.
  Both exhibit the pairing $\phi^{(1,4)}<\phi^{(0)}<\phi^{(2,3)}$
  with ratio $\phi^{(2,3)}/\phi^{(1,4)}\approx 2$.}
\label{tab:floors_both_L12}
\begin{tabular}{cllll}
\toprule
Class $k$ & \multicolumn{2}{c}{\mfld{m003} (Meyerhoff)} &
            \multicolumn{2}{c}{\mfld{m006}} \\
          & $\phi_{\rm floor}$ ($\degr$) & Word
          & $\phi_{\rm floor}$ ($\degr$) & Word \\
\midrule
0 & 0.04373 & \word{aaaabbbaabaB} & 0.02300 & \word{aaaaaaaaabAb} \\
1 & 0.04319 & \word{aaBAbAAbAABB} & 0.00530 & \word{aaaaaababAAA} \\
2 & 0.09387 & \word{aabAAbAAABB}  & 0.01060 & \word{aaababaaabab} \\
3 & 0.09387 & \word{aaaaBabbAAB}  & 0.01060 & \word{AAABABAAABAB} \\
4 & 0.04319 & \word{aabbAAbaBaaB} & 0.00530 & \word{aaaBAAABAAAA} \\
\midrule
\multicolumn{2}{l}{$\phi^{(2,3)}/\phi^{(1,4)}$} & $2.17$ &
  & $2.00$ \\
\bottomrule
\end{tabular}
\end{table}

\subsection{Spectral floors at length 12: the 4:2:1 hierarchy}
Extending the census to length~12, all five classes develop small floors
and the initial asymmetry evolves into a stable hierarchy.
Table~\ref{tab:floors_L12} shows the length-12 floors.

\begin{table}[tp]
\caption{Spectral floors per homology class at word length $\leq 12$
  for \mfld{m006} ($1{,}062{,}802$ genuine loxodromics).
  The $4:2:1$ ratio (0.023 : 0.011 : 0.005) is stable across lengths
  $\geq 10$.}
\label{tab:floors_L12}
\begin{tabular}{clll}
\toprule
Class $k$ & $\phi_{\rm floor}^{(k)}(12)$ ($\degr$) & Word & Length \\
\midrule
0 & 0.02300 & \word{aaaaaaaaabAb} & 12 \\
1 & 0.00530 & \word{aaaaaababAAA} & 12 \\
2 & 0.01060 & \word{aaababaaabab} & 12 \\
3 & 0.01060 & \word{AAABABAAABAB} & 12 \\
4 & 0.00530 & \word{aaaBAAABAAAA} & 12 \\
\midrule
\multicolumn{2}{l}{Ratio $\phi^{(0)}:\phi^{(2,3)}:\phi^{(1,4)}$}
& \multicolumn{2}{l}{$4.3:2:1$} \\
\bottomrule
\end{tabular}
\end{table}

\subsection{Exponential decay with class-dependent rates}
Figure~\ref{fig:decay} plots $\log\phi_{\rm floor}^{(k)}(L)$ versus $L$
for $L=1,\ldots,12$.
Over the range $L=1,\ldots,12$, each class exhibits approximately log-linear decay, consistent with
\begin{equation}
  \phi_{\rm floor}^{(k)}(L) \sim A_k\,e^{-c_k L}.
  \label{eq:exp_decay}
\end{equation}
Table~\ref{tab:rates} gives the fitted rates $c_k$ and AIC scores
comparing exponential against power-law fits.
The exponential model is preferred for all five classes.

\begin{table}[tp]
\caption{Exponential decay rates $c_k$ from log-linear fits to
  $\phi_{\rm floor}^{(k)}(L)$, with 95\% bootstrap confidence intervals
  (2000 resamples) and $\Delta$AIC $=$ AIC(exp) $-$ AIC(pow).
  Negative $\Delta$AIC favours exponential over power-law.
  The pairings $c_1=c_4$ and $c_2=c_3$ are confirmed:
  their CIs overlap completely.}
\label{tab:rates}
\begin{tabular}{cccr}
\toprule
Class $k$ & $c_k$ & 95\% CI & $\Delta$AIC \\
\midrule
0 & 0.750 & [0.434,\ 0.944] & $-33.6$ \\
1 & 0.985 & [0.541,\ 1.376] & $-49.0$ \\
2 & 0.401 & [0.080,\ 0.645] & $-4.0$  \\
3 & 0.401 & [0.079,\ 0.651] & $-4.0$  \\
4 & 0.985 & [0.585,\ 1.431] & $-49.0$ \\
\midrule
\multicolumn{4}{l}{$\Delta$AIC $=$ AIC(exp) $-$ AIC(pow); negative favours exponential.}\\
\multicolumn{4}{l}{CIs from $2000$ bootstrap resamples (word-length basis).}\\
\bottomrule
\end{tabular}
\end{table}

\begin{figure}[tp]
\includegraphics[width=0.85\textwidth]{fig_floor_decay.pdf}
\caption{Spectral floor $\phi_{\rm floor}^{(k)}(L)$ (log scale) vs.\ word length $L$ for \mfld{m006}, partitioned by homology class $k\in\mathbb{Z}/5$. Horizontal axis: word length $L=1,\ldots,12$. Vertical axis: minimum folded twist angle (degrees, log scale). Dashed lines: exponential fits $A_k e^{-c_k L}$ with rates from Table~\ref{tab:rates}. Classes~1 and~4 decay fastest ($c\approx 0.985$, red/orange); classes~2 and~3 intermediate ($c\approx 0.401$, green/purple); class~0 slowest ($c\approx 0.750$, blue). The $4:2:1$ floor ratio at $L=12$ is the principal result (Table~\ref{tab:floors_L12}).}
\label{fig:decay}
\end{figure}

The pairing $c_1=c_4$ and $c_2=c_3$ and the isolation of class~0 are
the principal structural observations.
Explicitly: classes $k=1$ and $k=4$ satisfy $1+4=5\equiv 0\pmod{5}$;
classes $k=2$ and $k=3$ satisfy $2+3=5\equiv 0\pmod{5}$;
class $k=0$ is self-paired.
This strongly suggests that the rate pairing reflects the involution
$k\mapsto -k$ (mod~5) on $\mathbb{Z}/5$.

\section{The Twist-Suppressed Geodesic}
\label{sec:aaabab}

The class-4 floor is set from length~6 onwards by the geodesic
\word{AAABAB} (and its class-1 inverse \word{aaabab}).
Its geometric invariants are:
\begin{align}
  |\lambda(\word{AAABAB})| &= 3.730, &
  \phi(\word{AAABAB}) &= 179.995\degr, &
  \phi_{\rm fold} &= 0.0053\degr, \notag\\
  \ell(\word{AAABAB}) &= 2.633, &
  \mathrm{Tr}(\word{AAABAB}) &= -3.998, &
  [{\rm AAABAB}] &= -6\equiv 4\pmod{5}. \notag
\end{align}
The eigenvalue $\lambda\approx -3.730$ is nearly real negative, so the
element acts as a near-pure translation of length $\ell=2.633$ with
$\phi\approx\pi$ ($\phi_{\rm fold}\approx 0$).
The trace $|\mathrm{Tr}|=3.998$ is far from $2$ (the identity value),
confirming this is a genuine geodesic, not a near-trivial element.
It belongs to homology class~4 (non-identity coset) despite having
nearly vanishing twist.

\section{Discussion}
\label{sec:discussion}

The main findings are:
\begin{enumerate}
\item All five homology classes of $\pi_1(\mfld{m006})$ exhibit
      exponential decay of their spectral floors with word length,
      with class-dependent rates $c_k$.

\item The rates form a $4:2:1$ hierarchy: $c_{1,4}\approx 0.985$,
      $c_{2,3}\approx 0.401$, $c_0\approx 0.750$, confirmed by AIC
      over lengths 1--12.

\item The pairing $c_1=c_4$, $c_2=c_3$ reflects the involution
      $k\mapsto -k$ (mod~5) on $\mathbb{Z}/5\cong H_1(\mfld{m006})$.
      Class~0 (identity coset) is self-paired and has an intermediate
      rate, not the fastest or slowest.

\item At length~12, the Meyerhoff manifold \mfld{m003} (also $H_1=\mathbb{Z}/5$)
      exhibits the \emph{same} pairing structure as \mfld{m006}:
      floors $\phi^{(1,4)}=0.04319\degr$, $\phi^{(2,3)}=0.09387\degr$,
      $\phi^{(0)}=0.04373\degr$, giving ratio
      $\phi^{(2,3)}/\phi^{(1,4)}\approx 2.17$.
      The pairing $c_1=c_4$, $c_2=c_3$ thus appears to be a property
      of the $\mathbb{Z}/5$ torsion structure shared by both manifolds,
      not specific to \mfld{m006}'s arithmetic.
\end{enumerate}

\subsection{Floors vanish: consequence of the Ambient PGT}
\label{sec:floors_zero}

\begin{proposition}
\label{prop:floors_zero}
Let $M$ be a compact hyperbolic 3-manifold with
$H_1(M,\mathbb{Z})=\mathbb{Z}/5$.
For each homology class $k\in\mathbb{Z}/5$,
$\phi_{\rm floor}^{(k)}(L)\to 0$ as $L\to\infty$.
\end{proposition}

\begin{proof}
By the Ambient Prime Geodesic Theorem~\cite{BalogBiggs23} and the
equidistribution of homology classes~\cite{Sunada85}, for any
$\epsilon>0$ the count of class-$k$ geodesics with $\phi_{\rm fold}<\epsilon$
and length $\leq T$ grows as $(4\epsilon/\pi)\cdot e^{2T}/(10T)\to\infty$.
Hence every class eventually contains geodesics with
$\phi_{\rm fold}<\epsilon$, so $\phi_{\rm floor}^{(k)}(L)\to 0$.
\end{proof}

\begin{remark}
Proposition~\ref{prop:floors_zero} shows convergence to zero but
says nothing about the rate. The class-dependent rates $c_k$
(Table~\ref{tab:rates}) are the new content of this paper.
\end{remark}

\subsection{Margulis asymptotics and the exponential decay law}
\label{sec:margulis}

Proposition~\ref{prop:floors_zero} guarantees $\phi_{\rm floor}^{(k)}(L)\to 0$
but does not predict the rate.  Here we explain why exponential decay
is the natural expectation from the prime geodesic theorem.

The prime geodesic theorem (Margulis~\cite{Margulis69}) gives
\begin{equation}
  \pi_\Gamma(T) \sim \frac{e^{2T}}{2T}, \qquad T\to\infty,
  \label{eq:pgt}
\end{equation}
for the count of primitive closed geodesics of length $\leq T$.
The Ambient PGT~\cite{BalogBiggs23} and the equidistribution of homology
classes~\cite{Sunada85} together give, for each class $k\in\mathbb{Z}/5$
and each arc $J_\epsilon = \{\theta : |\phi_{\rm fold}|<\epsilon\}$
of measure $|J_\epsilon|=4\epsilon$,
\begin{equation}
  N^{(k)}(T,\epsilon)
  \;:=\;\#\{\gamma\in C_k : \ell(\gamma)\leq T,\;\phi_{\rm fold}(\gamma)<\epsilon\}
  \;\sim\; \frac{4\epsilon}{\pi}\cdot\frac{1}{5}\cdot\frac{e^{2T}}{2T}.
  \label{eq:count}
\end{equation}
The spectral floor $\phi_{\rm floor}^{(k)}(L)$ is the value of $\epsilon$
at which $N^{(k)}(T_L,\epsilon)$ first reaches~1,
where $T_L$ is the maximum geodesic length achievable at word length $L$.
Since each generator has a positive translation length bounded below by
the systole $\mathrm{sys}(\mfld{m006})>0$,
the translation length of a word of length $L$ is bounded above by
$L\cdot\ell_{\rm max}$ and below by $\ell_{\rm min}$, giving
$T_L\leq L\cdot\ell_{\rm max}$; empirically $T_L\sim\alpha L$ with
$\alpha\approx\ell_{\rm max}(\mfld{m006})\approx 1.5$ for the word lengths studied.
Since $T_L \sim \alpha L$ for some constant $\alpha>0$ determined by the
generator lengths, equation~\eqref{eq:count} gives
\begin{equation}
  \phi_{\rm floor}^{(k)}(L)
  \;\sim\; C_k \cdot e^{-2\alpha L},
  \label{eq:floor_asymptotic}
\end{equation}
where the prefactor $C_k$ encodes the leading-order correction to the
equidistribution of twist angles within class~$k$.
The Margulis asymptotics predict that the \emph{leading rate} $2\alpha$
is the same for all classes, but the \emph{prefactor} $C_k$ can differ.
The class-dependent rates $c_k$ observed in our data (Table~\ref{tab:rates})
are therefore corrections to the leading term --- they reflect
class-to-class variation in the density of near-zero-twist geodesics
beyond what equidistribution predicts.

The pairing $c_1=c_4$, $c_2=c_3$ (Table~\ref{tab:rates},
Figures~\ref{fig:decay}) is then a statement about the symmetry of
these corrections under the involution $k\mapsto -k$ (mod~5).
We conjecture that this symmetry is arithmetic in origin.

\begin{conjecture}
The decay rates $c_k$ are determined by the arithmetic invariants of
$\pi_1(\mfld{m006})$.
Specifically, the pairing $c_k=c_{5-k}$ for $k\neq 0$ reflects an
outer involution $\rho\mapsto\bar{\rho}$ of the holonomy representation
$\rho:\pi_1(\mfld{m006})\to\PSL$, arising from complex conjugation in
the invariant trace field.
Manifolds with the same invariant trace field and quaternion algebra as
\mfld{m006} have the same decay rate hierarchy $(c_0,c_1,c_2,c_3,c_4)$.
\end{conjecture}

A proof of this conjecture would require computing the invariant trace
field of \mfld{m006} via the \textsc{Snap} package~\cite{Snap} and
relating the arithmetic involution to the observed rate pairing.
A direct test is feasible: \textsc{Snap} computes the invariant trace
field $k(M)$ and quaternion algebra $\mathcal{A}$ of any cusped
hyperbolic manifold.  Running \textsc{Snap} on the Dehn-filling
description of \mfld{m006} and checking whether $\mathcal{A}$ admits
a complex-conjugation involution would either confirm or refute the
arithmetic origin of the pairing $c_1=c_4$, $c_2=c_3$.
We note that \mfld{m003} is proven arithmetic (invariant trace field of
discriminant $-283$)~\cite{Chinburg98}; the arithmeticity of \mfld{m006}
is expected but unproven.

The exponential decay rates $c_k$ provide a new family of geometric
invariants of compact hyperbolic 3-manifolds with torsion in $H_1$.
Comparing these rates across the OrientableClosedCensus for other manifolds
with $H_1=\mathbb{Z}/p$ ($p$ prime) would test the conjecture and may
reveal which arithmetic properties control the rate hierarchy.

\begin{remark}[Scope of the decay claim]
The exponential fits in Table~\ref{tab:rates} and Figure~\ref{fig:decay}
are empirical descriptions of the data over the range $L=1,\ldots,12$.
No claim is made regarding the asymptotic behaviour of
$\phi_{\rm floor}^{(k)}(L)$ as $L\to\infty$ beyond the convergence to
zero established in Proposition~\ref{prop:floors_zero}.
The Margulis asymptotics (Section~\ref{sec:margulis}) suggest exponential
decay is natural, but a rigorous proof of the rate is an open problem.
\end{remark}

We note that the word-length $L$ depends on the choice of generators
$\{a,b\}$ for $\pi_1(\mfld{m006})$.  A generator-independent formulation
would use the geodesic translation length $\ell(\gamma)=2\log|\lambda(\gamma)|$
in place of $L$.  A preliminary analysis using the running maximum of
$\ell(\gamma)$ as the $x$-axis variable confirms that the class pairing
$c_1=c_4$ and $c_2=c_3$ persists, though the rate values differ;
a rigorous treatment via the geodesic length spectrum of \mfld{m006}
is deferred to future work.

\section*{Data Availability}
All census data and analysis scripts supporting this study are publicly
available at \url{https://github.com/drmlgentry/hyperbolic-flavor-scan}.
The complete word-length-12 census files (\texttt{twist\_census\_len12\_m003.csv}
and \texttt{twist\_census\_len12\_m006.csv}) are included in the repository.

\section*{Acknowledgements}
The author thanks the \textsc{SnapPy} development team.
Computations used \textsc{SnapPy}~\cite{SnapPy}, SciPy, NumPy, and Pandas.
All code and data are publicly available at
\url{https://github.com/drmlgentry/hyperbolic-flavor-scan}.

\section*{Conflict of Interest}
The author declares no conflicts of interest.

\begin{thebibliography}{99}
\bibitem{GentryCKM}
  M.~L.~Gentry,
  ``CKM Quark Mixing from Geodesic Axes of Compact Hyperbolic 3-Manifolds,''
  Phys.\ Rev.\ D, submitted March 2026 (DQ14014).
\bibitem{GentryPMNS}
  M.~L.~Gentry,
  ``PMNS Lepton Mixing from the Borel Structure of Hyperbolic Holonomy,''
  Phys.\ Rev.\ D, submitted March 2026 (DQ14050).
\bibitem{GentryTwist}
  M.~L.~Gentry,
  ``Twist Angle Spectrum of Hyperbolic Holonomy,''
  Phys.\ Rev.\ D, submitted March 2026 (es2026mar17\_1007).
\bibitem{BalogBiggs23}
  A.~Balog, A.~Biggs, A.~de~Faveri, A.~Harper, K.~Pushkarev, S.~Steiner,
  ``Ambient Prime Geodesic Theorems on Hyperbolic 3-Manifolds,''
  Int.\ Math.\ Res.\ Not.\ \textbf{2023}, no.~1, 588--606.
\bibitem{Margulis69}
  G.~A.~Margulis,
  ``Certain applications of ergodic theory to the investigation of manifolds
  of negative curvature,'' Funct.\ Anal.\ Appl.\ \textbf{3}, 335--336 (1969).
\bibitem{Sunada85}
  T.~Sunada,
  ``Riemannian coverings and isospectral manifolds,''
  Ann.\ Math.\ \textbf{121}, 169--186 (1985).
\bibitem{MaclachlanReid}
  C.~Maclachlan and A.~W.~Reid,
  \textit{The Arithmetic of Hyperbolic 3-Manifolds},
  Springer, New York (2003).
\bibitem{Meyerhoff87}
  R.~Meyerhoff,
  ``Hyperbolic 3-manifolds with equal volumes but different Chern-Simons
  invariants,''
  in \textit{Low-dimensional topology and Kleinian groups},
  London Math.\ Soc.\ Lecture Note Ser.~112 (1986), pp.~209--215.
\bibitem{Chinburg98}
  T.~Chinburg, E.~Friedman, K.~N.~Jones, A.~W.~Reid,
  ``The arithmetic hyperbolic 3-manifold of smallest volume,''
  Ann.\ Scuola Norm.\ Sup.\ Pisa \textbf{30}, 1 (2001).
\bibitem{Snap}
  D.~Coulson, O.~Goodman, C.~Hodgson, W.~Neumann,
  ``Computing arithmetic invariants of 3-manifolds,''
  Experimental Math.\ \textbf{9}, 127 (2000).
\bibitem{SnapPy}
  M.~Culler, N.~Dunfield, M.~Goerner, J.~Weeks,
  \textit{SnapPy}, \url{http://snappy.computop.org}.
\end{thebibliography}

\end{document}
